\chapter{Evaluating causal assistance without changing the question}
\label{la:evaluation}

The question throughout this chapter is whether assistance preserves a banking decision while improving its resolution at an acceptable cost. The opening defines the comparison; the coverage screen and prospective experiment then show what a finite evaluation could establish. The source studies test the reference answers, exposure history, adjudication and scoring rules on which that experiment would rely. Finally, the evaluator, memory and toolkit comparisons follow those requirements into a working interface. Their different failure mechanisms cannot be reduced to a single accuracy score.

\section{The benchmark unit is a causal problem}

An agent can produce a plausible report, execute valid code and still answer
the wrong causal question. Evaluation must follow the entire chain from
business meaning through assumptions and identification to uncertainty,
availability and decision loss. A benchmark of text fluency, graph edit
distance or successful tool calls measures only part of that chain.

CLadder demonstrates the value of formal questions with varied verbalizations,
including commonsense and deliberately unfamiliar or conflicting stories.
The inspected paper also identifies language-to-formal-tool formulation as
a continuing challenge. Its historical binary-answer accuracy is evidence
about those tasks and models; it is not a measurement of effect coverage,
decision loss or banking readiness. \citep{v12_la_2312_04350}

CausalDS expands the evaluation object to narrated synthetic mechanisms,
data and a separate observation model. Its inspected results show why
symbolic success, uncertainty quality, abstention and cost should be reported
separately. Its author-reported exam and repeated runs are useful design
evidence, not model evaluations conducted in this project. A generated world
with known truth also remains different from authentic bank data whose
semantic and assignment mechanisms must be established externally.
\citep{v12_la_2607_08093}

The proposed experimental unit is therefore an independently generated or
independently adjudicated causal mechanism and question. Several phrasings,
label conditions or stochastic model calls can be nested within that unit.
They do not create new independent worlds. Repeating one problem fifty times
can measure response variation for that problem; it does not provide fifty
independent demonstrations of general causal competence.

\section{Matched information and meaningful alternatives}

The prospective design compares a fixed numerical workflow, a language-only
workflow, retrieval assistance, a hybrid of numerical tools and typed model
proposals, and structured expert workflows with and without assistance.
Each arm needs a declared information budget. If one receives the true
causal graph and another must discover it, their effect-estimation scores
do not isolate the value of language assistance. If retrieval adds a valid
institutional fact, that information gain should be distinguished from a
better inference algorithm.

Eight strata are proposed: independent sufficient DAGs; latent-variable
graph classes; temporal data with known clocks; treatment-confounder
feedback; panel or quasi-experimental assignment; proxies, linkage and
selection; weak or nonidentifying effects; and ambiguous business targets
or monetary loss. Label and evidence conditions cross ordinary descriptions,
unfamiliar labels, misleading or irrelevant material, and conflict between
numerical evidence and text. These strata test different failure mechanisms
rather than treating all causal questions as interchangeable.

The current planning allocation is 400 mechanisms, with 50 per stratum and
three stochastic repetitions nested within each mechanism. This is an
allocation for developing the design, not a frozen final sample size or a
power calculation. For orientation, the worst-case normal approximation for
one Bernoulli proportion at 400 independent observations has half-width
\(1.96\sqrt{0.25/400}\), approximately 0.049. Paired differences, stratified
outcomes and rare severe errors require their own precision design. A
development pilot would inform that design and must be excluded from the
final test. The six task families developed in Section~\ref{exec13:live-design} describe workflow functions, whereas the eight strata here describe scientific conditions; they are different design dimensions. That later section uses additional sample sizes to illustrate error control, power and cost. Those calculations refine the choice of a trial rather than combine into an approved allocation.

Both natural-workflow and matched-budget comparisons can be useful. The
first asks what an analyst can achieve with a realistic workflow, including
review. The second asks what changes when information or computation is
controlled. Neither should hide extra expert time, repeated retrieval or
retry costs. The comparison must state which question it answers.

\section{A target error can survive perfect arithmetic}

An exact project IV construction illustrates a benchmark-label hazard.
Let a randomized binary instrument \(Z\) affect treatment for half the
population, the compliers; the other half never receives treatment. Set
\(Y(0)=0\) for everyone, \(Y(1)=1\) for compliers and \(Y(1)=3\) for
never-takers. Treatment equals \(Z\) for compliers and zero otherwise.
Randomization, exclusion and monotonicity hold in this construction.

The observed outcome difference across instrument values is \(0.5\) and
the treatment-probability difference is \(0.5\), so the Wald ratio is one.
The population ATE is two. Changing the never-takers' unobserved treated
outcome changes the population ATE without changing the observed law. A
scorer that calls the unrestricted Wald ratio the population ATE would
reward a target substitution. This is a finite counterexample supporting
review of generator assumptions, not a claim that every IV question in a
named benchmark uses this construction.

The same concern applies to trimming, overlap weighting and selected
populations. A language answer can reproduce the right formula for a
different target. Evaluation should therefore compare the population,
measure, action, comparator, horizon and units before comparing the numeric
answer. Critical semantic errors deserve their own count, even when a
numeric discrepancy happens to be small in the test mechanism.

\section{Graphs, sets and refusals need separate scores}

Graph evaluation should compare adjacency and endpoint statements against
the appropriate identifiable class. A justified circle or undirected edge
should not be penalized merely because a simulator internally selected one
fully directed graph. Orientation informativeness can be reported alongside
validity. This makes a method that preserves genuine ambiguity distinguishable
from one that returns little information and from one that confidently
overorients the graph.

Identification outcomes also need separate labels. Scientific
nonidentification, unsupported implementation, malformed output, explicit
abstention, timeout and missing input are not the same event. A well-designed
agent can correctly identify a missing premise without successfully computing
an effect. A backend can be unable to implement a valid functional. Combining
these into one ``incorrect'' or ``no answer'' category would obscure the
remedy.

For identified scalar targets, bias and error can be meaningful. For reported
sets, evaluate coverage, topology, finite availability and size. Whole-line
answers can have high coverage and little decision value. A method that
returns short intervals on a selected easy subset needs an all-request
availability denominator. Unbounded and disconnected sets must remain in the
evaluation rather than being replaced by finite plotting limits.

For an eligible finite interval \([l,u]\), one prospective score at nominal
level \(1-\alpha\) is
\[
 (u-l)+\frac{2}{\alpha}(l-\theta)1\{\theta<l\}
       +\frac{2}{\alpha}(\theta-u)1\{\theta>u\}.
\]
The formula balances length with penalties for missing the target. It is
not an instruction to compress a union of intervals into one interval. For
unions, coverage and total Lebesgue length can be reported while retaining
infinite lengths. Raw losses should not be capped simply to make a plot or
ranking convenient. The eligible target type must be stated before the
score is selected.

\section{Contamination, adaptation and selective reporting}

An exposed demonstration cannot become a held-out evaluation merely because
its variables are renamed. The local problem fixtures and replay examples
have already been used to develop and debug the prototype. They can continue
to test conformance, but they cannot provide independent estimates of a
language model's causal accuracy.

The final design needs a frozen development/test split, mechanism identities,
prompts, model/provider revisions, retrieval corpus and ranking rules, numerical
baselines, stopping rules and analysis code. Mechanism families should be
separated where family reuse would reveal the answer. Label novelty and
mechanism novelty should be varied independently. An unfamiliar name attached
to a memorized mechanism is not the same test as an unfamiliar mechanism.

Test outcomes cannot be used to tune prompt templates, choose the best
retriever, strengthen selected priors or discard difficult cells while
retaining a final-test interpretation. A changed method requires a new
reviewed evaluation revision. Failed and unavailable requests remain in the
record. Reporting only the critic-approved subset would confound answer
quality with selection by the critic.

The inspected OCI paper provides a useful refusal example: its critic
declines many replications when overlap is poor. That can be preferable
to an unsupported estimate, but it makes the denominator essential.
The reported successful answers, refused requests, review effort and
remaining target population must all be visible. The external results
are author reports, not a validation of this project's critic.
\citep{v12_la_2607_22443}

\section{Decision value and cost}

The ultimate question is whether assistance improves feasible decisions at
an acceptable total cost. In a known synthetic mechanism, a policy's expected
loss can be calculated under the mechanism and compared with the chosen
baseline. In a bank, the feasible policy, valuation model, risk limits,
operating costs and outcome truth need authentic inputs. Synthetic decision
loss does not establish institutional return on investment.

The prospective cost record includes model calls, input/output/cache tokens,
retrieval, numerical CPU time, wall time, memory, retries and human review
and rework. Actual dated prices are needed to turn these quantities into a
currency total. Unknown prices remain missing; zero live calls in an offline
prototype do not imply a free live workflow. A fast answer that requires
substantial correction can have a worse total cost than a slower transparent
analysis.

No final live-model experiment is reported by this chapter. Five prerequisite
classes remain: an exact live backend and custody terms; actual prices and
a bounded budget; independent held-out mechanisms and original benchmark
bytes; named business/diagram adjudicators and a rubric; and qualified
numerical baselines. These are missing program inputs. No credentials were
read and no provider endpoint was probed to infer that a service is unavailable.
The design provides a route to evaluation once those inputs are supplied;
it provides no current empirical model-performance claim. The remaining studies turn these requirements into inspectable tests: first the probability of passing a coverage rule, then a prospective comparison, and finally the reliability of its references, evaluators and numerical handoffs.
